\chapter{Исследовательская часть}
\section{Технические характеристики}

Технические характеристики устройства, на котором выполнялись замеры времени, представлены далее.

\begin{enumerate}
	\item Процессор	AMD Ryzen 5 5500u ЦПУ 2592 МГц, ядер: 8, логических процессоров: 12.
	\item Оперативная память: 16 ГБ.
	\item Операционная система: Fedora 39.
\end{enumerate}

При замерах времени ноутбук был включен в сеть электропитания. Запущены были только системные приложения.

\section{Демонстрация работы программы}

На рисунке~\ref{img:workExample} представлена демонстрация работы разработанного приложения для решения задачи коммивояжёра.

\includeimage
{workExample}
{f}
{H}
{0.75\textwidth}
{Пример работы программы}

\section{Временные характеристики}

Результаты замеров времени работы алгоритмов решения задачи коммивояжёра представлены в таблице~\ref{tbl:time} и на графике~\ref{plt:time}. Для каждого N проводилось 100 замеров времени, после чего результаты усреднялись. Все результаты замеров приведены в миллисекундах.

\begin{longtable}{|p{.33\textwidth - 2\tabcolsep}|p{.33\textwidth - 2\tabcolsep}|p{.33\textwidth - 2\tabcolsep}|}
	\caption{Результаты замеров времени работы алгоритмов}\label{tbl:time}  \\\hline
	Размерность матрицы N & Алгоритм решения полным перебором, мс & Алгоритм решения муравьиной колонией, мс   \\\hline
	\endfirsthead
	2 & 0.000001 & 0.000031 \\
	\hline
	3 & 0.000001 & 0.000103 \\
	\hline
	4 & 0.000002 & 0.000240 \\
	\hline
	5 & 0.000002 & 0.000480 \\
	\hline
	6 & 0.000009 & 0.000815 \\
	\hline
	7 & 0.000052 & 0.001332 \\
	\hline
	8 & 0.000351 & 0.002125 \\
	\hline
	9 & 0.002828 & 0.003229 \\
	\hline
	10 & 0.026926 & 0.004760 \\
	\hline
	11 & 0.190294 & 0.004048 \\
	\hline
\end{longtable}

\begin{figure}[H]
	\centering
	\includesvg[inkscapelatex=false, width=0.9\textwidth, height=0.9\textheight]{inc/img/graph.svg}
	\caption{Результаты замеров времени работы алгоритмов решения задачи коммивояжёра}
	\label{plt:time}
\end{figure}

\section*{Вывод}
На основе полученных был сделан вывод о том, что алгоритм решения задачи коммивояжёра полным перебором работает медленнее муравьиного алгоритма при большей размерности матрицы. Таким образом при размерности матрицы более 9 алгоритм муравьиной колонии по времени работает более эффективно.

\section{Проведение параметризации}

Автоматическая параметризация была проведена на трех классах данных.
Алгоритм был запущен для набора значений $\alpha, \beta, e  \in (0, 1)$. Количество дней всегда равно 50.

Итоговая таблица значений параметризации состоит из следующих колонок:
\begin{itemize}[label=---]
	\item $\alpha$~---~коэффициент жадности;
	\item $\beta$~---~коэффициент стадности;
	\item $e$~---~коэффициент испарения;
	\item \textit{Ошибка}~---~разность между полученным муравьиным алгоритмом решением и самым эффективным решением на рассматриваемых.
\end{itemize}

Цель параметризации~---~опредение комбинации параметров, которые позволяют решить задачу с наименьшим значением ошибки для выбранного класса данных.
Качество решения зависит от погрешности измерений.


\subsection{Класс данных 1}\label{par:class1}

Класс данных 1 представляет собой 3 матрицы смежности для 10 вершин~(небольшой разброс длины пути~---~от 1 до 10), каждая из которых представлена далее.

\begin{equation}
	\label{eq:kd1.1}
	K_{1.1} = \begin{pmatrix}
		0 & 9 & 9 & 2 & 2 & 4 & 9 & 6 & 8 & 5 \\ 
		9 & 0 & 1 & 8 & 6 & 4 & 10 & 4 & 1 & 10 \\ 
		9 & 1 & 0 & 10 & 5 & 1 & 5 & 1 & 3 & 8 \\ 
		2 & 8 & 10 & 0 & 5 & 6 & 9 & 5 & 1 & 9 \\ 
		2 & 6 & 5 & 5 & 0 & 8 & 1 & 9 & 9 & 5 \\ 
		4 & 4 & 1 & 6 & 8 & 0 & 2 & 7 & 10 & 9 \\ 
		9 & 10 & 5 & 9 & 1 & 2 & 0 & 3 & 10 & 7 \\ 
		6 & 4 & 1 & 5 & 9 & 7 & 3 & 0 & 10 & 6 \\ 
		8 & 1 & 3 & 1 & 9 & 10 & 10 & 10 & 0 & 8 \\ 
		5 & 10 & 8 & 9 & 5 & 9 & 7 & 6 & 8 & 0
	\end{pmatrix}
\end{equation}

\begin{equation}
	\label{eq:kd1.2}
	K_{1.2} = \begin{pmatrix}
		0 & 6 & 6 & 7 & 5 & 2 & 9 & 9 & 3 & 1 \\ 
		6 & 0 & 7 & 7 & 8 & 7 & 3 & 10 & 7 & 2 \\ 
		6 & 7 & 0 & 3 & 6 & 10 & 7 & 9 & 8 & 8 \\ 
		7 & 7 & 3 & 0 & 8 & 3 & 10 & 6 & 4 & 7 \\ 
		5 & 8 & 6 & 8 & 0 & 3 & 9 & 4 & 1 & 3 \\ 
		2 & 7 & 10 & 3 & 3 & 0 & 5 & 9 & 4 & 7 \\ 
		9 & 3 & 7 & 10 & 9 & 5 & 0 & 1 & 10 & 5 \\ 
		9 & 10 & 9 & 6 & 4 & 9 & 1 & 0 & 8 & 8 \\ 
		3 & 7 & 8 & 4 & 1 & 4 & 10 & 8 & 0 & 8 \\ 
		1 & 2 & 8 & 7 & 3 & 7 & 5 & 8 & 8 & 0  
	\end{pmatrix}
\end{equation}

\begin{equation}
	\label{eq:kd1.3}
	K_{1.3} = \begin{pmatrix}
		0 & 9 & 7 & 1 & 1 & 2 & 1 & 9 & 2 & 10 \\ 
		9 & 0 & 9 & 9 & 4 & 10 & 4 & 8 & 6 & 8 \\ 
		7 & 9 & 0 & 6 & 9 & 8 & 1 & 5 & 8 & 4 \\ 
		1 & 9 & 6 & 0 & 2 & 8 & 5 & 6 & 7 & 2 \\ 
		1 & 4 & 9 & 2 & 0 & 5 & 6 & 10 & 6 & 8 \\ 
		2 & 10 & 8 & 8 & 5 & 0 & 1 & 8 & 9 & 3 \\ 
		1 & 4 & 1 & 5 & 6 & 1 & 0 & 9 & 9 & 3 \\ 
		9 & 8 & 5 & 6 & 10 & 8 & 9 & 0 & 5 & 8 \\ 
		2 & 6 & 8 & 7 & 6 & 9 & 9 & 5 & 0 & 9 \\ 
		10 & 8 & 4 & 2 & 8 & 3 & 3 & 8 & 9 & 0
	\end{pmatrix}
\end{equation}

Для данного класса данных параметризация приведена в приложении~А.

Лучшими наборами этого класса являются:

\begin{enumerate}[label=---]
	\item $(\alpha = 0.1, \beta = 0.9, e = 0)$;
	\item $(\alpha = 0.1, \beta = 0.9, e = 0.5)$;
	\item $(\alpha = 0.1, \beta = 0.9, e = 0.6)$; 
	\item $(\alpha = 0.2, \beta = 0.8, e = 0.3)$; 
	\item $(\alpha = 0.2, \beta = 0.8, e = 0.4)$; 
	\item $(\alpha = 0.2, \beta = 0.8, e = 0.8)$; 
	\item $(\alpha = 0.2, \beta = 0.8, e = 0.9)$; 
	\item $(\alpha = 0.3, \beta = 0.7, e = 0.6)$;
	\item $(\alpha = 0.3, \beta = 0.7, e = 0.8)$; 
	\item $(\alpha = 0.4, \beta = 0.6, e = 0.7)$; 
	\item $(\alpha = 0.4, \beta = 0.6, e = 0.9)$.
\end{enumerate}

При этих значениях $\alpha$, $\beta$ и $e$ среднее значение ошибки равно 0, поэтому рекомендуется использовать именно их.

\subsection{Класс данных 2}\label{par:class2}

Класс данных 2 представляет собой 3 матрицы смежности для 10 вершин~(большой разброс значений длины пути~---~от 100 до 1100), каждая из которых представлена далее.



\begin{equation}
	\label{eq:kd2.1}
	K_{2.1} = \begin{pmatrix}
	0 & 295 & 769 & 471 & 188 & 880 & 293 & 107 & 996 & 234 \\ 
	295 & 0 & 372 & 287 & 622 & 695 & 474 & 476 & 976 & 663 \\ 
	769 & 372 & 0 & 528 & 477 & 968 & 106 & 1078 & 856 & 552 \\ 
	471 & 287 & 528 & 0 & 672 & 923 & 367 & 895 & 364 & 328 \\ 
	188 & 622 & 477 & 672 & 0 & 820 & 279 & 997 & 910 & 1085 \\ 
	880 & 695 & 968 & 923 & 820 & 0 & 498 & 104 & 1092 & 395 \\ 
	293 & 474 & 106 & 367 & 279 & 498 & 0 & 956 & 1084 & 301 \\ 
	107 & 476 & 1078 & 895 & 997 & 104 & 956 & 0 & 480 & 680 \\ 
	996 & 976 & 856 & 364 & 910 & 1092 & 1084 & 480 & 0 & 394 \\ 
	234 & 663 & 552 & 328 & 1085 & 395 & 301 & 680 & 394 & 0  
	\end{pmatrix}
\end{equation}

\begin{equation}
	\label{eq:kd2.2}
	K_{2.2} = \begin{pmatrix}
	0 & 856 & 276 & 676 & 1003 & 653 & 546 & 1010 & 351 & 1021 \\ 
	856 & 0 & 182 & 923 & 846 & 169 & 439 & 829 & 397 & 878 \\ 
	276 & 182 & 0 & 1008 & 1014 & 408 & 995 & 413 & 412 & 988 \\ 
	676 & 923 & 1008 & 0 & 708 & 270 & 692 & 628 & 369 & 273 \\ 
	1003 & 846 & 1014 & 708 & 0 & 922 & 126 & 168 & 499 & 748 \\ 
	653 & 169 & 408 & 270 & 922 & 0 & 441 & 664 & 378 & 692 \\ 
	546 & 439 & 995 & 692 & 126 & 441 & 0 & 587 & 461 & 236 \\ 
	1010 & 829 & 413 & 628 & 168 & 664 & 587 & 0 & 334 & 530 \\ 
	351 & 397 & 412 & 369 & 499 & 378 & 461 & 334 & 0 & 294 \\ 
	1021 & 878 & 988 & 273 & 748 & 692 & 236 & 530 & 294 & 0  
	\end{pmatrix}
\end{equation}

\begin{equation}
	\label{eq:kd2.3}
	K_{2.3} = \begin{pmatrix}
	0 & 1063 & 546 & 1072 & 692 & 461 & 381 & 588 & 774 & 412 \\ 
	1063 & 0 & 196 & 102 & 301 & 508 & 630 & 570 & 681 & 453 \\ 
	546 & 196 & 0 & 315 & 469 & 572 & 964 & 810 & 137 & 961 \\ 
	1072 & 102 & 315 & 0 & 403 & 624 & 1041 & 539 & 577 & 191 \\ 
	692 & 301 & 469 & 403 & 0 & 733 & 261 & 638 & 425 & 853 \\ 
	461 & 508 & 572 & 624 & 733 & 0 & 999 & 425 & 1060 & 675 \\ 
	381 & 630 & 964 & 1041 & 261 & 999 & 0 & 738 & 875 & 677 \\ 
	588 & 570 & 810 & 539 & 638 & 425 & 738 & 0 & 939 & 284 \\ 
	774 & 681 & 137 & 577 & 425 & 1060 & 875 & 939 & 0 & 927 \\ 
	412 & 453 & 961 & 191 & 853 & 675 & 677 & 284 & 927 & 0  
	\end{pmatrix}
\end{equation}

Для данного класса данных параметризация приведена в приложении~A.

Лучшим набором этого класса является $(\alpha = 0, \beta = 1, e = 0.8)$.
При этих значениях $\alpha$, $\beta$ и $e$ среднее значение ошибки минимально исходя из полученных данных, поэтому рекомендуется использовать именно их.

